% BÀI TOÁN HÌNH HỌC - BAI_HINH_39 (TRÍCH ĐỀ 39)
\begin{ex}
% Câu 12
	Cho $\triangle ABC$ có 3 góc nhọn và đường cao $BE$. Gọi $H, K$ lần lượt là chân đường vuông góc kẻ từ điểm $E$ đến $AB, BC$.
	\begin{enumerate}[a)]
		\item Chứng minh tứ giác $BHEK$ nội tiếp;
		\item Chứng minh: $BH \cdot BA = BK \cdot BC$; Gọi $F$ là chân đường vuông góc kẻ từ điểm $C$ đến đường thẳng $AB$, $I$ là trung điểm của đoạn thẳng $EF$. Chứng minh rằng $H,I,K$ thẳng hàng.
	\end{enumerate}
\loigiai{
{
}
\end{ex}
